\documentclass[UTF8,a4paper,11pt,openany]{ctexbook}
\usepackage{../common/statlearnbook}
\SLSetBookMetadata{第 1 册}{统计学习方法讲义 v2.6.0}{数学基础与统计学习基本理论}
\begin{document}
\setcounter{page}{98}
\setcounter{chapter}{6}
\setcounter{figure}{0}
\pagestyle{headings}
\chaptermark{信息论基础}
\begin{theorem}[有限均匀分布最大熵]
若$X$有$K$个可能值，则 $0\le H(X)\le\log K$；上界当且仅当分布均匀时取得。
\end{theorem}
\paragraph{核验路线。}先写证明目标与合法条件，标出第一个可执行数学动作；下方严格证明保留全部技术细节。
\begin{proof}
记正概率支持为$S=\{x:p(x)>0\}$。由$\log$的凹性与Jensen不等式，
\[
H(P)=\sum_{x\in S}p(x)\log\frac1{p(x)}
\le\log\!\left(\sum_{x\in S}p(x)\frac1{p(x)}\right)
=\log|S|\le\log K.
\]
上界等号当且仅当$P$在全部$K$点上均匀；每个$-p(x)\log p(x)$非负，故$H(P)\ge0$。
\end{proof}
% v2.3.1 figure UID: FIG-P092-01
% Proposed post-v2.3.1 polish for FIG-P092-01: protect key-label size and stroke hierarchy.
\tikzset{
  slfig-FIG-P092-01/.style={every path/.append style={line cap=round,line join=round}},
  slfig-FIG-P092-01-guide/.style={draw=SLGray!76,line width=.52pt,
    dash pattern=on 3pt off 2pt},
  slfig-FIG-P092-01-label/.style={font=\fontsize{9.2pt}{11.2pt}\selectfont,
    fill=white,fill opacity=.92,text opacity=1,inner sep=1.2pt}
}
\pgfplotsset{slfig-FIG-P092-01-axis/.style={
  tick label style={font=\fontsize{8.8pt}{10.6pt}\selectfont},
  label style={font=\fontsize{9.4pt}{11.4pt}\selectfont},clip=false}}
\pgfkeysifdefined{/tikz/alt}{}{\tikzset{alt/.code={}}}
\begin{figure}[H]
\centering
\begin{tikzpicture}[slfig-FIG-P092-01,alt={对象—关系—结论：二元熵在p=1/2处达到1比特，在确定分布两端趋于0}]
\begin{axis}[slfig-FIG-P092-01-axis,
  slfig axis,width=.78\linewidth,height=5.8cm,
  axis lines=left,xlabel={$p$},ylabel={$H_2(p)$},
  xmin=-.04,xmax=1.04,ymin=-.04,ymax=1.15,
  xtick={0,.5,1},xticklabels={$0$,$\tfrac12$,$1$},
  ytick={0,1},yticklabels={$0$,$1$},
  domain=.001:.999,samples=321,clip=false
]
  \addplot[draw=SLBlue,line width=1.05pt,no marks]
    {-(x*ln(x)+(1-x)*ln(1-x))/ln(2)};
  \addplot[only marks,mark=*,mark size=2.15pt,draw=SLBlue,fill=SLBlue]
    coordinates {(0,0) (.5,1) (1,0)};
  \draw[slfig-FIG-P092-01-guide] (axis cs:.5,0)--(axis cs:.5,1);
  \draw[slfig-FIG-P092-01-guide] (axis cs:0,1)--(axis cs:.5,1);
  \node[slfig-FIG-P092-01-label,anchor=south]
    at (axis cs:.5,1.025) {最大不确定性：$1$ 比特};
  \node[slfig-FIG-P092-01-label,anchor=base west,text=SLGray!90]
    at (axis cs:.02,.055) {确定性};
  \node[slfig-FIG-P092-01-label,anchor=base east,text=SLGray!90]
    at (axis cs:.98,.055) {确定性};
  \node[slfig direct label,slfig-FIG-P092-01-label,fill=none,anchor=center]
    at (axis cs:.65,.45)
    {$H_2(p)=H_2(1-p)$};
\end{axis}
\end{tikzpicture}
\caption{二元熵在$p=1/2$处达到1比特，在确定分布两端趋于0}\label{fig:V1-C06-binary-entropy}
\end{figure}

对称性给出$H_{\mathrm b}(p)=H_{\mathrm b}(1-p)$；二阶导数为负说明\cref{fig:V1-C06-binary-entropy}中的曲线严格凹，驻点$p=1/2$因此是唯一最大点。
\SLReviewSection{条件熵与联合熵}{sec:v260-fig092-context}
\end{document}
